\begin{pastexam}[label={px:ncku-113-ee-01}]{成大電機所 113 學年度 是非題第 1 題}
（是非；每題 3 分，答「錯」要指出原因，答錯倒扣 1 分）

\begin{tikzpicture}[baseline=(a.base),>={Latex[length=4pt]},v/.style={draw,circle,inner sep=1.2pt,font=\small}]
  \node[v] (a) at (0,0) {2}; \node[v] (b) at (1,0) {3};
  \draw[->] ([yshift=2pt]a.east)--([yshift=2pt]b.west);
  \draw[->] ([yshift=-2pt]b.west)--([yshift=-2pt]a.east);
\end{tikzpicture} is a strongly connect component of the graph
\begin{tikzpicture}[baseline=(a.base),>={Latex[length=4pt]},v/.style={draw,circle,inner sep=1.2pt,font=\small}]
  \node[v] (a) at (0,0) {1}; \node[v] (b) at (1,0) {2}; \node[v] (c) at (2,0) {3};
  \node[v] (d) at (3,0) {5}; \node[v] (e) at (4,0) {4};
  \foreach \s/\t in {a/b,b/c,c/d} {
    \draw[->] ([yshift=2pt]\s.east)--([yshift=2pt]\t.west);
    \draw[->] ([yshift=-2pt]\t.west)--([yshift=-2pt]\s.east);
  }
  \draw[->] (d)--(e);
\end{tikzpicture}.
\end{pastexam}

\begin{pastsoln}
這份考卷和成大電通所 113 學年度的「資料結構」是同一份（電機系、電通所、智慧資安學程共用），完整詳解見歷屆試題~\ref{px:ncku-113-cce-01}。

錯。 這張圖的強連通元件是 $\{1, 2, 3, 5\}$ 和 $\{4\}$；$\{2, 3\}$ 雖然兩兩互相可達，但不是極大的，不是強連通元件。
\end{pastsoln}
